\section{Conclusion}
\label{sec:conclusion}

This paper presented \mra{}, a configurable benchmark layer that isolates vehicle-specific and simulator-specific behavior behind a replaceable adapter; the reported implementation uses HoloOcean and a BlueROV2-class agent. It defines single-vehicle evaluation, then extends it to staggered cooperative fleets with team scoring. Controllers see only the official observation, while an independent referee uses privileged state for crossings, penalties, timeouts, ranking and scores. We formalized the benchmark instance, gate geometry, single-rover and team scoring and staggered fleet semantics, then grounded them in the implementation.

In HoloOcean clean validation, the deterministic rule baseline completes all three official tracks in each of five seeds under the official observation contract. The staggered two-rover smoke demonstration also finishes in every seed with no inter-vehicle proximity events. Under currents, the same baseline degrades sharply: the \texttt{medium} profile finishes in only $3/5$ seeds and the \texttt{strong} profile finishes in none. These failures remain part of the benchmark evidence because gates, margins and current profiles are unchanged.

In the matched HoloOcean comparison on clean Horseshoe Bay, both camera-assisted rule controllers finish every seed. Replacing continuous visual correction near the aperture with center-then-commit passage reduces mean official time from $231.2$ to $184.9$\,s and mean gate/world collisions from $0.6$ to zero. A separate coordination diagnostic uses heterogeneous beacon-only navigation baselines with three rovers, no currents or obstacles, two start gaps and three seeds. All six coordinated runs finish with no inter-vehicle proximity events or gate/world collisions. Uncoordinated teams incur many gate/world collisions and sometimes fail to finish.
